\documentclass{beamer}
\usetheme{Madrid}
\usepackage[utf8]{inputenc}
\usepackage{amsmath}
\usepackage{amsfonts}
\usepackage{amssymb}
\usepackage{graphicx}
\usepackage{xcolor}
\usepackage{tikz}

\title{CAMS Exam Preparation}
\subtitle{Domain B: Global AML/CFT Frameworks \& Regulations (100 Study Slides)}
\author{Based on CAMS Candidate Handbook \& Candidate Feedback}
\date{}

\setbeamercovered{transparent}
\setbeamertemplate{navigation symbols}{}

\begin{document}
	
	% TITLE SLIDE
	\begin{frame}
		\titlepage
		\centering
		\textit{Domain B: 20\% of CAMS Exam | 100 Key Concepts \& Explanations}
	\end{frame}
	
	% ============================================================
	% SECTION B.1: FATF - OVERVIEW AND MUTUAL EVALUATIONS
	% ============================================================
	
	\begin{frame}{B.1.1: What is FATF and What Does It Do?}
		\begin{block}{FATF Overview}
			\textbf{Answer:} FATF is an inter-governmental body established in 1989 by the G7. It sets global AML/CFT standards through the 40 Recommendations. The Secretariat is at OECD in Paris. FATF has 39 members and monitors implementation through Mutual Evaluations. It also identifies high-risk jurisdictions (Grey List and Black List).
		\end{block}
		\textcolor{blue}{Key point:} FATF Recommendations are not legally binding treaties, but countries must implement them to avoid being blacklisted and facing economic consequences.
	\end{frame}
	
	\begin{frame}{B.1.2: The Structure of the 40 FATF Recommendations}
		\begin{block}{Seven Categories}
			\textbf{Answer:} The 40 Recommendations are organized into seven categories:
			\begin{enumerate}
				\item AML/CFT Policies and Coordination (Rec. 1-2)
				\item Money Laundering and Confiscation (Rec. 3-4)
				\item Terrorist Financing and Proliferation Financing (Rec. 5-8)
				\item Preventive Measures (Rec. 9-23)
				\item Transparency of Legal Persons/Arrangements (Rec. 24-25)
				\item Powers of Competent Authorities (Rec. 26-35)
				\item International Cooperation (Rec. 36-40)
			\end{enumerate}
		\end{block}
	\end{frame}
	
	\begin{frame}{B.1.3: How Does FATF Assess Countries?}
		\begin{block}{Mutual Evaluation - Two Pillars}
			\textbf{Answer:} FATF assesses countries on two pillars:
			\begin{itemize}
				\item \textbf{Technical Compliance}: Do laws exist on paper? Rated Compliant (C), Largely Compliant (LC), Partially Compliant (PC), or Non-Compliant (NC).
				\item \textbf{Effectiveness}: Do laws achieve intended outcomes? Rated High, Substantial, Moderate, or Low across 11 Immediate Outcomes (IO.1-IO.11).
			\end{itemize}
			\textcolor{red}{Critical:} A country can have perfect laws (high Technical Compliance) but Low Effectiveness if no prosecutions occur.
		\end{block}
	\end{frame}
	
	\begin{frame}{B.1.4: What Are the 11 FATF Immediate Outcomes?}
		\begin{block}{IO.1 through IO.11}
			\textbf{Answer:} The 11 Immediate Outcomes measure effectiveness:
			\begin{itemize}
				\item \textbf{IO.1}: ML/TF risks understood and coordinated
				\item \textbf{IO.2}: International cooperation used
				\item \textbf{IO.3}: Supervision of FIs and DNFBPs
				\item \textbf{IO.4}: Preventive measures applied
				\item \textbf{IO.5}: Legal persons prevented from misuse
				\item \textbf{IO.6}: Financial intelligence used by authorities
				\item \textbf{IO.7}: ML investigated and prosecuted
				\item \textbf{IO.8}: Proceeds of crime confiscated
				\item \textbf{IO.9}: TF investigated and prosecuted
				\item \textbf{IO.10}: TF asset freezing and confiscation
				\item \textbf{IO.11}: Counter-proliferation financing measures
			\end{itemize}
		\end{block}
	\end{frame}
	
	\begin{frame}{B.1.5: What Does Low Effectiveness on IO.6 Look Like?}
		\begin{block}{IO.6 - Financial Intelligence Used}
			\textbf{Answer:} IO.6 assesses whether financial intelligence from the FIU is used by law enforcement.
			\vspace{0.2cm}
			\textbf{Low effectiveness example:} The FIU receives 60,000 STRs annually but never disseminates any intelligence to law enforcement. Despite high volume, the intelligence is not used, so IO.6 is rated Low.
			\vspace{0.2cm}
			\textcolor{blue}{Required evidence:} Dissemination of intelligence, investigations initiated from FIU reports, and feedback from law enforcement to FIU.
		\end{block}
	\end{frame}
	
	\begin{frame}{B.1.6: What Does Low Effectiveness on IO.7 Look Like?}
		\begin{block}{IO.7 - ML Investigated and Prosecuted}
			\textbf{Answer:} IO.7 assesses whether money laundering is investigated and prosecuted.
			\vspace{0.2cm}
			\textbf{Low effectiveness example:} High STR volume but only three ML investigations in five years and zero convictions. Prosecutors lack training on financial crime methodology. This results in Low effectiveness regardless of STR volume.
			\vspace{0.2cm}
			\textcolor{blue}{Required evidence:} ML investigations initiated from FIU intelligence, actual convictions, trained prosecutors.
		\end{block}
	\end{frame}
	
	\begin{frame}{B.1.7: What Does Low Effectiveness on IO.8 Look Like?}
		\begin{block}{IO.8 - Proceeds Confiscated}
			\textbf{Answer:} IO.8 assesses whether proceeds of crime are confiscated.
			\vspace{0.2cm}
			\textbf{Low effectiveness example:} Confiscation laws exist but no assets ever confiscated because prosecutors cannot trace criminal proceeds and no asset management office exists. Laws alone are insufficient - actual confiscation is required.
			\vspace{0.2cm}
			\textcolor{blue}{Required evidence:} Assets identified, traced, frozen, and confiscated; asset management office; non-conviction based forfeiture available.
		\end{block}
	\end{frame}
	
	\begin{frame}{B.1.8: What Does FATF Recommendation 1 Require?}
		\begin{block}{Rec. 1 - Risk-Based Approach}
			\textbf{Answer:} Countries and financial institutions must identify, assess, and understand ML/TF risks, then apply measures commensurate with those risks.
			\vspace{0.2cm}
			\textbf{Key principles:}
			\begin{itemize}
				\item Simplified measures where risks are lower (Simplified Due Diligence - SDD)
				\item Enhanced measures where risks are higher (Enhanced Due Diligence - EDD)
				\item Risk assessments must be updated regularly
			\end{itemize}
			\textcolor{blue}{Not allowed:} One-size-fits-all approach or blanket de-risking without individual assessment.
		\end{block}
	\end{frame}
	
	\begin{frame}{B.1.9: What Does FATF Recommendation 2 Require?}
		\begin{block}{Rec. 2 - National Cooperation}
			\textbf{Answer:} Countries must establish national AML/CFT policies, designate competent authorities (FIU, supervisors, law enforcement), and ensure cooperation and coordination among them.
			\vspace{0.2cm}
			\textbf{Deficiency example:} A financial supervisor identifies significant AML deficiencies at banks but never communicates these findings to the FIU or law enforcement. This violates Recommendation 2.
			\vspace{0.2cm}
			\textcolor{blue}{Required:} Information sharing mechanisms, regular coordination meetings, adequate resources.
		\end{block}
	\end{frame}
	
	\begin{frame}{B.1.10: What Does FATF Recommendation 3 Require?}
		\begin{block}{Rec. 3 - Money Laundering Offense}
			\textbf{Answer:} Countries must criminalize money laundering based on the widest range of predicate offenses.
			\vspace{0.2cm}
			\textbf{Key requirements:}
			\begin{itemize}
				\item Apply to all serious predicate offenses (the "all-crimes approach" is preferred)
				\item Self-laundering MUST be criminalized (predicate offender can be charged with ML)
				\item Apply regardless of where the predicate offense occurred
				\item Minimum predicate offenses include: drug trafficking, fraud, corruption, tax crimes
			\end{itemize}
			\textcolor{blue}{Example:} Drug trafficker who processes own proceeds can be charged with both drug trafficking AND money laundering.
		\end{block}
	\end{frame}
	
	\begin{frame}{B.1.11: What Does FATF Recommendation 4 Require?}
		\begin{block}{Rec. 4 - Confiscation}
			\textbf{Answer:} Countries must adopt measures to confiscate proceeds of crime and instrumentalities of crime.
			\vspace{0.2cm}
			\textbf{Key requirements:}
			\begin{itemize}
				\item Include property of corresponding value (substitute assets)
				\item Extend to property converted or transformed from proceeds
				\item Include income or benefits derived from proceeds
				\item Powers to identify, trace, freeze, and seize property
			\end{itemize}
			\textcolor{blue}{Important:} Non-conviction based forfeiture (civil forfeiture) is recommended.
		\end{block}
	\end{frame}
	
	\begin{frame}{B.1.12: What Does FATF Recommendation 5 Require?}
		\begin{block}{Rec. 5 - Terrorist Financing Offense}
			\textbf{Answer:} Countries must criminalize terrorist financing as a standalone offense.
			\vspace{0.2cm}
			\textbf{Key requirements:}
			\begin{itemize}
				\item Criminalize TF regardless of whether funds are actually used for a terrorist act
				\item \textbf{No minimum monetary threshold} (even \$1 is sufficient)
				\item Collecting or providing funds for terrorist purposes is enough
				\item Apply to nationals and entities
			\end{itemize}
			\textcolor{red}{Deficiency example:} A \$10,000 threshold for TF violates Recommendation 5 because it implies small amounts are legal.
		\end{block}
	\end{frame}
	
	\begin{frame}{B.1.13: What Does FATF Recommendation 6 Require?}
		\begin{block}{Rec. 6 - Targeted Financial Sanctions (TF)}
			\textbf{Answer:} Countries must implement UN Security Council resolutions on terrorist financing.
			\vspace{0.2cm}
			\textbf{Key requirements:}
			\begin{itemize}
				\item Implement UNSCR 1267, 1373, and 1988
				\item Freeze assets without delay of designated persons
				\item Prohibit funds or assets from being available to designated persons
				\item Establish mechanisms for delisting and access to funds for basic expenses
			\end{itemize}
			\textcolor{red}{Freeze must be immediate - no prior notice to the customer.}
		\end{block}
	\end{frame}
	
	\begin{frame}{B.1.14: What Does FATF Recommendation 7 Require?}
		\begin{block}{Rec. 7 - Proliferation Financing}
			\textbf{Answer:} Countries must implement UN Security Council resolutions on WMD proliferation.
			\vspace{0.2cm}
			\textbf{Key requirements:}
			\begin{itemize}
				\item Implement UNSCR 1540, 1718, 1737, 1747, 1803, 1929, 2231, 2270
				\item Freeze assets of designated proliferators
				\item Prohibit provision of funds or services to designated persons
				\item Apply enhanced due diligence to proliferation financing risks
			\end{itemize}
			\textcolor{blue}{Red flags:} Dual-use goods exports to sanctioned jurisdictions, payment routing through multiple intermediaries.
		\end{block}
	\end{frame}
	
	\begin{frame}{B.1.15: What Does FATF Recommendation 8 Require?}
		\begin{block}{Rec. 8 - Non-Profit Organizations}
			\textbf{Answer:} Countries must apply targeted risk-based supervision of NPOs vulnerable to terrorist financing abuse.
			\vspace{0.2cm}
			\textbf{Key requirements:}
			\begin{itemize}
				\item Mechanisms to ensure NPOs are not misused by terrorist organizations
				\item Outreach and awareness programs for NPOs
				\item Monitoring of cross-border transactions involving NPOs
			\end{itemize}
			\textcolor{blue}{Red flags for NPOs:} Large donation from a Grey List jurisdiction, rapid cash withdrawal (80\%+ within 48 hours) far from headquarters.
		\end{block}
	\end{frame}
	
	\begin{frame}{B.1.16: What Does FATF Recommendation 9 Require?}
		\begin{block}{Rec. 9 - Financial Institution Secrecy Laws}
			\textbf{Answer:} Financial institution secrecy laws must NOT impede implementation of FATF Recommendations.
			\vspace{0.2cm}
			\textbf{Key requirements:}
			\begin{itemize}
				\item FIs must provide customer information to competent authorities
				\item FIs cannot invoke bank secrecy to refuse STR filing
				\item Professional secrecy (attorney-client privilege) has limited exceptions for financial transactions
			\end{itemize}
			\textcolor{red}{Bank secrecy is NOT a valid defense against AML obligations.}
		\end{block}
	\end{frame}
	
	\begin{frame}{B.1.17: When Must CDD Be Conducted Under FATF Rec. 10?}
		\begin{block}{Rec. 10 - CDD Triggers}
			\textbf{Answer:} CDD must be conducted in four specific situations:
			\begin{enumerate}
				\item When establishing a business relationship
				\item When carrying out an occasional transaction above the threshold (USD/EUR 15,000)
				\item When there is a suspicion of money laundering or terrorist financing (regardless of amount)
				\item When there is doubt about the veracity of previously obtained identification documents
			\end{enumerate}
			\textcolor{blue}{Note:} CDD is required even for very small transactions if suspicion exists.
		\end{block}
	\end{frame}
	
	\begin{frame}{B.1.18: What Are the Four Required Elements of CDD?}
		\begin{block}{Rec. 10 - CDD Elements}
			\textbf{Answer:} CDD must include four elements:
			\begin{enumerate}
				\item Identify the customer and verify their identity using reliable, independent source documents
				\item Identify the beneficial owner and take reasonable measures to verify their identity
				\item Understand the purpose and intended nature of the business relationship
				\item Conduct ongoing due diligence on the business relationship and scrutinize transactions throughout the relationship
			\end{enumerate}
			\textcolor{red}{All four elements are mandatory.}
		\end{block}
	\end{frame}
	
	\begin{frame}{B.1.19: What Is the Recordkeeping Requirement Under FATF Rec. 11?}
		\begin{block}{Rec. 11 - Recordkeeping}
			\textbf{Answer:} All records must be retained for a minimum of 5 years.
			\vspace{0.2cm}
			\textbf{Specific requirements:}
			\begin{itemize}
				\item Transaction records (national and international): 5 years following transaction completion
				\item Customer identification records (CDD files): 5 years after business relationship ends
				\item STRs and supporting documentation: 5 years after filing
				\item Records must be available to competent authorities upon request
			\end{itemize}
			\textcolor{red}{Minimum 5 years} - some countries require longer (e.g., 10 years in some EU states).
		\end{block}
	\end{frame}
	
	\begin{frame}{B.1.20: What CDD Obligations Apply to Casinos Under FATF Rec. 12?}
		\begin{block}{Rec. 12 - DNFBPs (Casinos)}
			\textbf{Answer:} Casinos must conduct CDD when a customer purchases chips or tokens equivalent to USD/EUR 3,000 or more.
			\vspace{0.2cm}
			\textbf{Additional requirements:}
			\begin{itemize}
				\item CDD includes verifying identity and source of funds for large transactions
				\item Casinos must be licensed or registered
				\item Casinos must be supervised for AML/CFT compliance
			\end{itemize}
			\textcolor{blue}{Note:} Some jurisdictions (e.g., U.S.) require CTRs at $10,000 and SARs at $5,000 threshold in addition to FATF's \$3,000 CDD trigger.
		\end{block}
	\end{frame}
	
	\begin{frame}{B.1.21: What CDD Obligations Apply to Real Estate Agents Under FATF Rec. 13?}
		\begin{block}{Rec. 13 - DNFBPs (Real Estate)}
			\textbf{Answer:} Real estate agents must conduct CDD when involved in transactions concerning buying or selling real estate.
			\vspace{0.2cm}
			\textbf{Key requirements:}
			\begin{itemize}
				\item Applies to both residential and commercial property
				\item Must verify identity of buyer and seller
				\item Must report suspicious transactions
				\item Recordkeeping for 5 years
			\end{itemize}
			\textcolor{blue}{Red flags:} All-cash purchases through anonymous LLCs, wire transfers from attorney IOLTA accounts, refusal to disclose beneficial owner.
		\end{block}
	\end{frame}
	
	\begin{frame}{B.1.22: What Are the MVTS Requirements Under FATF Rec. 14?}
		\begin{block}{Rec. 14 - Money or Value Transfer Services}
			\textbf{Answer:} All money or value transfer service providers must be licensed or registered and subject to AML/CFT monitoring.
			\vspace{0.2cm}
			\textbf{Key requirements:}
			\begin{itemize}
				\item CDD, recordkeeping, and STR filing obligations apply
				\item Unlicensed MVTS providers subject to criminal or administrative sanctions
				\item \textbf{No cultural heritage exemption} for hawala operators
			\end{itemize}
			\textcolor{red}{Hawala operators must be licensed like any other MVTS provider. "Traditional" or "cultural" is not an exemption.}
		\end{block}
	\end{frame}
	
	\begin{frame}{B.1.23: What Does FATF Rec. 15 Require for Virtual Assets?}
		\begin{block}{Rec. 15 - New Technologies}
			\textbf{Answer:} Countries must assess ML/TF risks of new technologies and regulate virtual asset service providers (VASPs).
			\vspace{0.2cm}
			\textbf{Key requirements:}
			\begin{itemize}
				\item VASPs must be licensed or registered
				\item VASPs subject to AML/CFT obligations (CDD, Travel Rule, SAR filing)
				\item \textbf{No exemption} for DeFi if control exists
				\item Countries must regulate virtual asset activities
			\end{itemize}
			\textcolor{red}{FATF updated Rec. 15 to explicitly include virtual assets.}
		\end{block}
	\end{frame}
	
	\begin{frame}{B.1.24: What Information Must Accompany a Wire Transfer Under FATF Rec. 16?}
		\begin{block}{Rec. 16 - Wire Transfer Rule (Travel Rule)}
			\textbf{Answer:} For cross-border wire transfers above USD/EUR 1,000, the following information is required:
			\vspace{0.2cm}
			\textbf{Originator information:}
			\begin{itemize}
				\item Full name
				\item Account number (or unique transaction reference)
				\item Address, national ID number, customer ID number, OR date and place of birth (at least ONE)
			\end{itemize}
			\textbf{Beneficiary information:}
			\begin{itemize}
				\item Full name
				\item Account number (or unique identifier)
			\end{itemize}
			\textcolor{red}{Receiving institution must obtain and hold this information.}
		\end{block}
	\end{frame}
	
	\begin{frame}{B.1.25: Can a Bank Rely on a Third Party for CDD Under FATF Rec. 17?}
		\begin{block}{Rec. 17 - Third-Party Reliance}
			\textbf{Answer:} Yes, but with strict conditions and the bank remains ultimately responsible.
			\vspace{0.2cm}
			\textbf{Conditions:}
			\begin{itemize}
				\item Third party must be subject to FATF-compliant AML requirements
				\item Ultimate responsibility for CDD remains with the bank (cannot outsource liability)
				\item Bank must have direct access to CDD information immediately
				\item Third party must provide copies of identification documents upon request
			\end{itemize}
			\textcolor{red}{Cannot outsource responsibility - the bank is always accountable.}
		\end{block}
	\end{frame}
	
	\begin{frame}{B.1.26: What Does FATF Rec. 18 Require for Financial Groups?}
		\begin{block}{Rec. 18 - Group-Wide Programs}
			\textbf{Answer:} Financial groups must implement group-wide AML/CFT programs across all branches and subsidiaries.
			\vspace{0.2cm}
			\textbf{Key requirements:}
			\begin{itemize}
				\item Apply to all branches and subsidiaries (including foreign)
				\item Group policies must be at least as strict as home country requirements
				\item Information sharing within the group is permitted (subject to confidentiality)
				\item Enhanced group measures for higher-risk jurisdictions
			\end{itemize}
			\textcolor{blue}{Example:} An EU parent bank must apply EU AML standards to its Cayman Islands subsidiary.
		\end{block}
	\end{frame}
	
	\begin{frame}{B.1.27: How Must a Bank Respond to FATF Grey List and Black List Jurisdictions?}
		\begin{block}{Rec. 19 - High-Risk Jurisdictions}
			\textbf{Answer:} Different responses apply:
			\vspace{0.2cm}
			\textbf{Black List (Call for Action - Iran, North Korea):}
			\begin{itemize}
				\item Apply enhanced due diligence (EDD)
				\item Apply countermeasures (restrict or terminate correspondent relationships)
			\end{itemize}
			\textbf{Grey List (Increased Monitoring - Turkey, Jamaica):}
			\begin{itemize}
				\item Apply EDD (enhanced scrutiny)
				\item Countermeasures not automatically required (risk-based)
				\item Individual risk assessments required (no blanket termination)
			\end{itemize}
		\end{block}
	\end{frame}
	
	\begin{frame}{B.1.28: When Must an STR Be Filed Under FATF Rec. 20?}
		\begin{block}{Rec. 20 - Suspicious Transaction Reporting}
			\textbf{Answer:} STR must be filed immediately upon suspicion, regardless of transaction value.
			\vspace{0.2cm}
			\textbf{Key principles:}
			\begin{itemize}
				\item File immediately upon suspicion (no delay to gather more evidence)
				\item No minimum monetary threshold
				\item No requirement for conclusive proof - reasonable suspicion is the standard
				\item No disclosure to customer (tipping-off prohibition)
				\item Safe harbor from civil and criminal liability for good-faith filings
			\end{itemize}
			\textcolor{red}{Reasonable suspicion is the standard - not certainty, not probable cause.}
		\end{block}
	\end{frame}
	
	\begin{frame}{B.1.29: What Is the Tipping-Off Prohibition Under FATF Rec. 21?}
		\begin{block}{Rec. 21 - Tipping-Off}
			\textbf{Answer:} FIs and their directors, officers, and employees are prohibited from disclosing that a SAR has been filed.
			\vspace{0.2cm}
			\textbf{Prohibited disclosures include:}
			\begin{itemize}
				\item Telling a customer a SAR has been filed on their account
				\item Disclosing that "compliance is reviewing your account" (implies SAR)
				\item Disclosing investigation to unauthorized third parties (press, public)
			\end{itemize}
			\textbf{Permissible disclosures:}
			\begin{itemize}
				\item Internal sharing with legal department on need-to-know basis
				\item 314(b) sharing with other opted-in institutions
				\item Responding to law enforcement subpoena
			\end{itemize}
		\end{block}
	\end{frame}
	
	\begin{frame}{B.1.30: What CDD Obligations Apply to Lawyers Under FATF Rec. 22?}
		\begin{block}{Rec. 22 - DNFBPs (Lawyers, Notaries)}
			\textbf{Answer:} Lawyers and notaries must conduct CDD when performing specified financial activities:
			\vspace{0.2cm}
			\textbf{Activities requiring CDD:}
			\begin{itemize}
				\item Buying or selling real estate for a client
				\item Managing client money, securities, or accounts
				\item Managing bank or savings accounts
				\item Creating or operating legal persons or trusts
			\end{itemize}
			\textbf{Exemption:} Providing legal advice on a client's legal position or representing a client in legal proceedings (CDD not required, but STR still required if suspicious).
		\end{block}
	\end{frame}
	
	\begin{frame}{B.1.31: What Obligations Apply to Trust and Company Service Providers Under FATF Rec. 23?}
		\begin{block}{Rec. 23 - DNFBPs (TCSPs)}
			\textbf{Answer:} Trust and Company Service Providers must conduct CDD and maintain records.
			\vspace{0.2cm}
			\textbf{Key requirements:}
			\begin{itemize}
				\item Conduct CDD when forming legal persons or arrangements
				\item Identify beneficial owners of the legal persons they create
				\item Maintain records for 5+ years
				\item File STRs when suspicious activity is detected
			\end{itemize}
			\textbf{Examples of TCSPs:} Company formation agents, registered agents, nominee directors and shareholders, trust administrators.
		\end{block}
	\end{frame}
	
	\begin{frame}{B.1.32: What Does FATF Rec. 24 Require for Beneficial Ownership of Legal Persons?}
		\begin{block}{Rec. 24 - Legal Person Transparency}
			\textbf{Answer:} Countries must ensure adequate, accurate, and current beneficial ownership information is obtained and held.
			\vspace{0.2cm}
			\textbf{Key requirements:}
			\begin{itemize}
				\item Prohibit bearer shares OR require immobilization with a regulated custodian
				\item BO information must be accessible to competent authorities
				\item Sanctions for non-compliance
			\end{itemize}
			\textbf{Deficiency example:} A BO registry exists but is not accessible to financial institutions without a court order (delay of weeks). This violates Rec. 24.
		\end{block}
	\end{frame}
	
	\begin{frame}{B.1.33: What Does FATF Rec. 25 Require for Beneficial Ownership of Trusts?}
		\begin{block}{Rec. 25 - Legal Arrangement Transparency}
			\textbf{Answer:} Trustees must hold adequate, accurate, and current beneficial ownership information on trusts.
			\vspace{0.2cm}
			\textbf{Who must be identified:}
			\begin{itemize}
				\item Settlor (the person who creates the trust)
				\item Trustee(s) (and UBOs if corporate trustee)
				\item Protector (if the protector has veto rights or power to replace the trustee)
				\item Beneficiaries (to the extent they are determinable)
			\end{itemize}
			\textcolor{red}{A protector with veto power over distributions IS a beneficial owner.}
		\end{block}
	\end{frame}
	
	\begin{frame}{B.1.34: What Does FATF Rec. 26 Require for Supervision of Financial Institutions?}
		\begin{block}{Rec. 26 - Supervision of FIs}
			\textbf{Answer:} Supervisors must have adequate powers and use a risk-based approach.
			\vspace{0.2cm}
			\textbf{Key requirements:}
			\begin{itemize}
				\item Supervisors must have powers to inspect and impose sanctions
				\item Use risk-based approach to allocate resources and determine examination frequency
				\item On-site AND off-site monitoring required
				\item Penalties for violations must be imposed (not optional)
			\end{itemize}
			\textbf{Deficiency example:} Same 5-year exam frequency for all banks regardless of risk; no penalties issued in 10 years.
		\end{block}
	\end{frame}
	
	\begin{frame}{B.1.35: What Powers Must Supervisors Have Under FATF Rec. 27?}
		\begin{block}{Rec. 27 - Powers of Supervisors}
			\textbf{Answer:} Supervisors must have comprehensive enforcement powers.
			\vspace{0.2cm}
			\textbf{Required powers:}
			\begin{itemize}
				\item Conduct inspections (on-site and off-site)
				\item Request access to all information (CDD files, STRs, transaction data)
				\item Impose sanctions (fines, restrictions, license revocation)
				\item Issue cease-and-desist orders
				\item Remove or suspend personnel
			\end{itemize}
			\textcolor{blue}{Supervisors must have "teeth" - power to enforce, not just recommend.}
		\end{block}
	\end{frame}
	
	\begin{frame}{B.1.36: What Is the Role of a Financial Intelligence Unit (FIU) Under FATF Rec. 29?}
		\begin{block}{Rec. 29 - Financial Intelligence Units}
			\textbf{Answer:} The FIU is the national center for receiving, analyzing, and disseminating financial intelligence.
			\vspace{0.2cm}
			\textbf{Key functions:}
			\begin{itemize}
				\item Receive and analyze STRs from FIs and DNFBPs
				\item Disseminate financial intelligence to competent authorities (law enforcement, prosecutors)
				\item Have access to necessary information (public databases, financial records)
				\item Secure exchange of information with foreign FIUs (via Egmont Group)
				\item Operational independence and autonomy
			\end{itemize}
			\textbf{Examples:} FinCEN (U.S.), UKFIU (UK), TRACFIN (France), UIF (Spain).
		\end{block}
	\end{frame}
	
	\begin{frame}{B.1.37: What Does FATF Rec. 30 Require for Law Enforcement?}
		\begin{block}{Rec. 30 - Law Enforcement Authorities}
			\textbf{Answer:} Law enforcement must conduct financial investigations into ML/TF.
			\vspace{0.2cm}
			\textbf{Key requirements:}
			\begin{itemize}
				\item LE must conduct parallel financial investigations in predicate crime cases
				\item Must be able to identify, trace, freeze, and seize assets
				\item Must have power to obtain financial records expeditiously (days, not months)
				\item Specialized investigators and prosecutors for financial crime
			\end{itemize}
			\textbf{Deficiency example:} 6-12 month delay to obtain bank records (assets can be moved by then).
		\end{block}
	\end{frame}
	
	\begin{frame}{B.1.38: What Investigative Powers Are Required Under FATF Rec. 31?}
		\begin{block}{Rec. 31 - Powers of Competent Authorities}
			\textbf{Answer:} Competent authorities must have comprehensive investigative powers.
			\vspace{0.2cm}
			\textbf{Required powers:}
			\begin{itemize}
				\item Compel production of documents and records (including financial records)
				\item Conduct searches of persons and premises
				\item Take witness statements under oath
				\item Obtain financial and law enforcement information
				\item Use special investigative techniques (undercover operations, electronic surveillance, controlled deliveries)
			\end{itemize}
			These powers apply to both ML and TF investigations.
		\end{block}
	\end{frame}
	
	\begin{frame}{B.1.39: What Controls Are Required for Cash Couriers Under FATF Rec. 32?}
		\begin{block}{Rec. 32 - Cash Couriers}
			\textbf{Answer:} Countries must have declaration or disclosure systems for cross-border currency movement.
			\vspace{0.2cm}
			\textbf{Required controls:}
			\begin{itemize}
				\item Declaration OR disclosure system for currency and bearer negotiable instruments at borders
				\item Authority to stop, restrain, and question cash couriers
				\item Sanctions for false declarations or non-disclosure (seizure, fines, criminal charges)
				\item Ability to seize cash suspected of being ML/TF related
			\end{itemize}
			\textbf{Deficiency example:} No declaration system, cash couriers intercepted with \$500,000 but no penalties.
		\end{block}
	\end{frame}
	
	\begin{frame}{B.1.40: What Statistics Must Countries Maintain Under FATF Rec. 33?}
		\begin{block}{Rec. 33 - Statistics}
			\textbf{Answer:} Countries must maintain comprehensive statistics on AML/CFT effectiveness.
			\vspace{0.2cm}
			\textbf{Data that must be collected:}
			\begin{itemize}
				\item STRs received, analyzed, and disseminated by the FIU
				\item ML/TF investigations, prosecutions, and convictions
				\item Asset freezing, seizure, and confiscation amounts
				\item International cooperation requests (MLATs sent, received, and executed)
				\item Mutual evaluation follow-up
			\end{itemize}
			Statistics are used to inform policy and resource allocation decisions.
		\end{block}
	\end{frame}
	
	\begin{frame}{B.1.41: What Guidance and Feedback Must Competent Authorities Provide Under FATF Rec. 34?}
		\begin{block}{Rec. 34 - Guidance and Feedback}
			\textbf{Answer:} Competent authorities must provide guidance and feedback to FIs and DNFBPs.
			\vspace{0.2cm}
			\textbf{What must be provided:}
			\begin{itemize}
				\item Red flag indicators for detecting suspicious activity
				\item Typologies and emerging ML/TF trends
				\item Feedback on STR quality and usefulness
				\item Best practices and case studies
				\item Outreach and awareness programs
			\end{itemize}
			\textcolor{red}{This is NOT optional - Rec. 34 is a standard, not a best practice.}
		\end{block}
	\end{frame}
	
	\begin{frame}{B.1.42: What Sanctions Are Required Under FATF Rec. 35?}
		\begin{block}{Rec. 35 - Sanctions}
			\textbf{Answer:} Sanctions for non-compliance must be effective, proportionate, and dissuasive.
			\vspace{0.2cm}
			\textbf{Types of sanctions:}
			\begin{itemize}
				\item Criminal sanctions (imprisonment, criminal fines)
				\item Civil sanctions (civil penalties, damages)
				\item Administrative sanctions (fines, license revocation, debarment, cease-and-desist orders)
				\item Available for FIs, DNFBPs, directors, and senior management
				\item Available for both natural persons (individuals) and legal persons (corporations)
			\end{itemize}
		\end{block}
	\end{frame}
	
	\begin{frame}{B.1.43: What Is the Purpose of Mutual Legal Assistance Treaties (MLATs) Under FATF Rec. 36?}
		\begin{block}{Rec. 36 - Mutual Legal Assistance Treaties}
			\textbf{Answer:} MLATs provide formal, legally binding mechanisms for cross-border exchange of evidence.
			\vspace{0.2cm}
			\textbf{Key requirements:}
			\begin{itemize}
				\item Provide the widest range of mutual legal assistance
				\item Respond to MLAT requests expeditiously
				\item Designate a central authority for MLAT processing
				\item Execute requests for evidence, witness testimony, restraint, and confiscation
			\end{itemize}
			\textcolor{blue}{MLAT vs. Egmont:} MLAT for evidence admissible in court (formal, slow); Egmont for intelligence sharing (informal, fast, not court-admissible).
		\end{block}
	\end{frame}
	
	\begin{frame}{B.1.44: How Does FATF Rec. 37 Address Dual Criminality?}
		\begin{block}{Rec. 37 - Dual Criminality}
			\textbf{Answer:} MLAT should be provided regardless of dual criminality requirements.
			\vspace{0.2cm}
			\textbf{Key principles:}
			\begin{itemize}
				\item MLAT assistance should be provided even if dual criminality is not satisfied
				\item If dual criminality is required, the predicate offense list should be interpreted broadly
				\item Treat money laundering as a predicate offense for terrorist financing (and vice versa)
				\item Avoid restrictive interpretations that hinder international cooperation
			\end{itemize}
			\textcolor{blue}{Goal:} Minimize barriers to international cooperation on ML/TF cases.
		\end{block}
	\end{frame}
	
	\begin{frame}{B.1.45: What Asset Recovery Assistance Is Required Under FATF Rec. 38?}
		\begin{block}{Rec. 38 - Asset Freezing and Confiscation}
			\textbf{Answer:} Countries must provide assistance for freezing, seizure, and confiscation of assets.
			\vspace{0.2cm}
			\textbf{Key requirements:}
			\begin{itemize}
				\item Provide assistance for freezing, seizure, and confiscation of proceeds of crime
				\item Non-conviction based forfeiture assistance
				\item Respond to requests for property of corresponding value (substitute assets)
				\item Manage and dispose of confiscated property
				\item Consider asset sharing agreements (splitting recovered assets with assisting country)
			\end{itemize}
		\end{block}
	\end{frame}
	
	\begin{frame}{B.1.46: What Extradition Requirements Does FATF Rec. 39 Impose?}
		\begin{block}{Rec. 39 - Extradition}
			\textbf{Answer:} Countries must extradite persons charged with ML/TF offenses.
			\vspace{0.2cm}
			\textbf{Key requirements:}
			\begin{itemize}
				\item Extradite persons charged with ML/TF offenses
				\item Simplified extradition procedures when possible
				\item Consider ML/TF as extraditable offenses in all extradition treaties
				\item Extradite own nationals where feasible (or prosecute them domestically)
			\end{itemize}
			\textcolor{blue}{Purpose:} Prevent criminals from evading justice by crossing borders.
		\end{block}
	\end{frame}
	
	\begin{frame}{B.1.47: What Other Forms of International Cooperation Are Required Under FATF Rec. 40?}
		\begin{block}{Rec. 40 - Other Forms of Cooperation}
			\textbf{Answer:} Countries must provide informal international cooperation beyond MLATs.
			\vspace{0.2cm}
			\textbf{Types of cooperation:}
			\begin{itemize}
				\item FIU-to-FIU cooperation (through the Egmont Group)
				\item Supervisor-to-supervisor cooperation (banking, securities, insurance regulators)
				\item Law enforcement-to-law enforcement cooperation (police-to-police)
				\item Provide spontaneous information sharing (without a request)
				\item No treaty required for informal cooperation
			\end{itemize}
		\end{block}
	\end{frame}
	
	% ============================================================
	% SECTION B.2: FATF GREY LIST AND BLACK LIST
	% ============================================================
	
	\begin{frame}{B.2.1: What Is the FATF Grey List and How Should Banks Respond?}
		\begin{block}{Grey List - Increased Monitoring}
			\textbf{Answer:} The Grey List includes jurisdictions with strategic AML/CFT deficiencies that have committed to an action plan.
			\vspace{0.2cm}
			\textbf{Bank obligations:}
			\begin{itemize}
				\item Apply Enhanced Due Diligence (EDD) to all transactions
				\item Increased scrutiny on correspondent accounts
				\item Individual risk assessments required (no blanket termination)
				\item NOT automatically required to terminate relationships
			\end{itemize}
			\textbf{Recent examples:} Turkey, Jamaica, Albania, Bulgaria (2023-2024).
		\end{block}
	\end{frame}
	
	\begin{frame}{B.2.2: What Is the FATF Black List and How Should Banks Respond?}
		\begin{block}{Black List - Call for Action}
			\textbf{Answer:} The Black List includes high-risk jurisdictions subject to a Call for Action.
			\vspace{0.2cm}
			\textbf{Bank obligations:}
			\begin{itemize}
				\item Apply countermeasures as called for by FATF
				\item Countermeasures may include: termination of correspondent relationships, restrictions on transactions, enhanced reporting requirements
				\item FATF calls on members to apply countermeasures
			\end{itemize}
			\textbf{Current Black List:} Iran, North Korea, Myanmar (as of 2024).
			\textcolor{red}{Black List requires stronger action than Grey List.}
		\end{block}
	\end{frame}
	
	% ============================================================
	% SECTION B.3: FATF-STYLE REGIONAL BODIES (FSRBs)
	% ============================================================
	
	\begin{frame}{B.3.1: How Do FSRBs Differ from FATF?}
		\begin{block}{FATF vs. FSRBs}
			\textbf{Answer:} FSRBs are autonomous regional bodies, not subordinate to FATF.
			\vspace{0.2cm}
			\textbf{Key distinctions:}
			\begin{itemize}
				\item FATF is a global body that sets the 40 Recommendations
				\item FSRBs are autonomous regional bodies (APG, MONEYVAL, GAFILAT, MENAFATF, GIABA, ESAAMLG)
				\item FSRBs promote, implement, and assess FATF Recommendations within their regions
				\item FSRBs conduct mutual evaluations using the same methodology as FATF
				\item Evaluation results are reported to FATF; significant deficiencies may lead to grey/black listing
			\end{itemize}
		\end{block}
	\end{frame}
	
	\begin{frame}{B.3.2: Can an FSRB Evaluate a Non-FATF Member?}
		\begin{block}{FSRB Evaluation Authority}
			\textbf{Answer:} Yes, FSRBs evaluate their member jurisdictions, whether or not those members are also FATF members.
			\vspace{0.2cm}
			\textbf{Key points:}
			\begin{itemize}
				\item FSRBs conduct mutual evaluations using the same FATF Methodology (Technical Compliance and Effectiveness across 11 Immediate Outcomes)
				\item Evaluation results are reported to FATF
				\item A country with significant deficiencies may be referred to FATF for potential public identification (Grey List or Black List)
			\end{itemize}
			\textcolor{red}{FSRB evaluation can lead to FATF grey listing - same consequences.}
		\end{block}
	\end{frame}
	
	% ============================================================
	% SECTION B.4: USA PATRIOT ACT
	% ============================================================
	
	\begin{frame}{B.4.1: What Is the USA PATRIOT Act and Why Is It Important for AML?}
		\begin{block}{USA PATRIOT Act Overview}
			\textbf{Answer:} The Uniting and Strengthening America by Providing Appropriate Tools Required to Intercept and Obstruct Terrorism Act was enacted after 9/11.
			\vspace{0.2cm}
			\textbf{Key features:}
			\begin{itemize}
				\item Expanded U.S. government surveillance and investigative powers
				\item Strengthened AML/CFT requirements for financial institutions
				\item Key sections 311-319 address correspondent banking and cross-border issues
				\item Administered by FinCEN, OFAC, and DOJ
			\end{itemize}
			\textcolor{red}{Heavily tested on CAMS - know Sections 311, 313, 314(a), 314(b), 319(a), 319(b).}
		\end{block}
	\end{frame}
	
	\begin{frame}{B.4.2: What Actions Can Be Taken Under USA PATRIOT Act Section 311?}
		\begin{block}{Section 311 - Primary Money Laundering Concern}
			\textbf{Answer:} The Secretary of Treasury can designate a foreign financial institution as a "primary money laundering concern."
			\vspace{0.2cm}
			\textbf{Measures that can be imposed:}
			\begin{itemize}
				\item Prohibit U.S. FIs from opening or maintaining correspondent or payable-through accounts for the designated institution
				\item Impose special recordkeeping and reporting requirements on transactions
				\item Impose conditions on correspondent accounts (e.g., requiring pre-approval of all transactions)
			\end{itemize}
			\textbf{Not authorized:} Seizing personal assets of employees without separate legal authority.
		\end{block}
	\end{frame}
	
	\begin{frame}{B.4.3: What Is the Shell Bank Prohibition Under USA PATRIOT Act Section 313?}
		\begin{block}{Section 313 - Shell Bank Prohibition}
			\textbf{Answer:} A "shell bank" is defined as a bank with no physical presence in any country and not affiliated with a regulated financial institution.
			\vspace{0.2cm}
			\textcolor{red}{\textbf{U.S. banks are absolutely prohibited}} from maintaining correspondent accounts for foreign shell banks.
			\vspace{0.2cm}
			\textbf{No exceptions for:}
			\begin{itemize}
				\item G20 central bank regulation (not relevant)
				\item Filing SARs every 90 days (does not cure)
				\item Enhanced monitoring (not sufficient)
			\end{itemize}
		\end{block}
	\end{frame}
	
	\begin{frame}{B.4.4: What Are a Bank's Obligations Under USA PATRIOT Act Section 314(a)?}
		\begin{block}{Section 314(a) - Law Enforcement Requests}
			\textbf{Answer:} FinCEN sends requests to financial institutions identifying subjects of investigation.
			\vspace{0.2cm}
			\textbf{Bank obligations:}
			\begin{itemize}
				\item Search accounts and transactions according to FinCEN's instructions
				\item Search for the specified time period (e.g., last 6 months, 1 year)
				\item Maintain confidentiality of the request (cannot disclose to customer)
				\item Report matches to law enforcement via FinCEN
			\end{itemize}
			\textbf{Not required:} Freeze accounts automatically (unless separate authority).
		\end{block}
	\end{frame}
	
	\begin{frame}{B.4.5: What Is Voluntary Information Sharing Under USA PATRIOT Act Section 314(b)?}
		\begin{block}{Section 314(b) - Voluntary Information Sharing}
			\textbf{Answer:} Financial institutions that have filed the required FinCEN opt-in notification may voluntarily share information with other opted-in institutions.
			\vspace{0.2cm}
			\textbf{Key features:}
			\begin{itemize}
				\item Share information about persons suspected of ML/TF activity
				\item Receive safe harbor from liability for sharing
				\item Build a complete picture across institutions
				\item May share SAR information (this is an exception to tipping-off rules)
			\end{itemize}
			\textcolor{blue}{314(b) sharing is explicitly permitted by law and is NOT tipping-off.}
		\end{block}
	\end{frame}
	
	\begin{frame}{B.4.6: What Is the Difference Between USA PATRIOT Act Sections 319(a) and 319(b)?}
		\begin{block}{Section 319(a) vs. 319(b)}
			\textbf{Answer:} The two sections address different aspects of correspondent banking.
			\vspace{0.2cm}
			\textbf{Section 319(a) - Forfeiture of Funds:}
			\begin{itemize}
				\item Allows U.S. government to seize funds from a correspondent account
				\item Funds must be traceable to money laundering
				\item Subject to judicial process
			\end{itemize}
			\textbf{Section 319(b) - Record Production:}
			\begin{itemize}
				\item Allows U.S. government to demand that a foreign bank produce records of its U.S. correspondent account
				\item If the foreign bank refuses, the U.S. bank may be required to terminate the account and funds may be seized
			\end{itemize}
		\end{block}
	\end{frame}
	
	% ============================================================
	% SECTION B.5: EU ANTI-MONEY LAUNDERING DIRECTIVES
	% ============================================================
	
	\begin{frame}{B.5.1: What Did the EU Fourth AML Directive (4AMLD) Introduce?}
		\begin{block}{4AMLD - Key Provisions (2015)}
			\textbf{Answer:} 4AMLD introduced several key requirements:
			\vspace{0.2cm}
			\begin{itemize}
				\item Beneficial ownership registers for corporate entities
				\item Risk-based approach to AML/CFT supervision
				\item Enhanced due diligence for PEPs (both domestic and foreign)
				\item Extended AML requirements to casinos and some virtual currency activities
				\item List of beneficial owners accessible to competent authorities
			\end{itemize}
			\textcolor{blue}{Note:} VASPs were fully added in 5AMLD, not 4AMLD.
		\end{block}
	\end{frame}
	
	\begin{frame}{B.5.2: What Did the EU Fifth AML Directive (5AMLD) Add?}
		\begin{block}{5AMLD - Key Provisions (2018)}
			\textbf{Answer:} 5AMLD expanded and strengthened AML requirements:
			\vspace{0.2cm}
			\begin{itemize}
				\item Extended CDD to virtual currency exchange platforms and custodian wallet providers (VASPs)
				\item Publicly accessible central registries of beneficial owners
				\item Lowered prepaid card CDD threshold from EUR 250 to EUR 150
				\item Bank account registers for FIUs to identify account holders
				\item Enhanced powers for FIUs (access to all financial information)
			\end{itemize}
			\textcolor{red}{5AMLD added VASPs to AML/CFT requirements.}
		\end{block}
	\end{frame}
	
	\begin{frame}{B.5.3: What Did the EU Sixth AML Directive (6AMLD) Introduce?}
		\begin{block}{6AMLD - Key Provisions (2020)}
			\textbf{Answer:} 6AMLD harmonized and strengthened criminal law provisions:
			\vspace{0.2cm}
			\begin{itemize}
				\item Harmonized 22 predicate offenses across all EU member states (binding, not optional)
				\item Extended criminal liability to legal persons (corporations) - not just fines
				\item Minimum maximum sentences: 4 years for ML, 5-15 years for aggravating circumstances
				\item No requirement for conviction of the predicate offender first
				\item Extended liability to negligent supervisors who fail to prevent ML
			\end{itemize}
			\textcolor{red}{6AMLD is heavily tested - know the differences from 5AMLD.}
		\end{block}
	\end{frame}
	
	% ============================================================
	% SECTION B.6: UN SANCTIONS FRAMEWORK
	% ============================================================
	
	\begin{frame}{B.6.1: What Are the Key UN Security Council Sanctions Regimes?}
		\begin{block}{UNSC Sanctions - Binding Obligations}
			\textbf{Answer:} UNSC resolutions are binding on all UN member states under international law.
			\vspace{0.2cm}
			\textbf{Key sanctions regimes:}
			\begin{itemize}
				\item UNSCR 1267/1988: Al-Qaida and Taliban
				\item UNSCR 1373: Terrorist financing (general)
				\item UNSCR 1540: WMD proliferation financing
				\item UNSCR 1718: North Korea (nuclear and missile programs)
				\item UNSCR 2231: Iran (nuclear program)
			\end{itemize}
			\textbf{Bank obligations:} Freeze assets immediately without prior notice, report to UNSC committee.
		\end{block}
	\end{frame}
	
	\begin{frame}{B.6.2: What Are Bank Obligations Under UNSCR 1540 for Proliferation Financing?}
		\begin{block}{UNSCR 1540 - WMD Non-Proliferation}
			\textbf{Answer:} UNSCR 1540 requires member states to establish controls to prevent WMD proliferation financing.
			\vspace{0.2cm}
			\textbf{Bank obligations:}
			\begin{itemize}
				\item Screen transactions for dual-use goods exports to sanctioned jurisdictions
				\item Identify proliferation financing red flags (inconsistent end-use, routing through intermediaries)
				\item File STRs for suspicious proliferation-related transactions
			\end{itemize}
			\textbf{Red flags:} Dual-use goods (radiation-hardened chips, maraging steel) to sanctioned countries, end-use claims inconsistent with goods (e.g., "atmospheric research" for missile components).
		\end{block}
	\end{frame}
	
	% ============================================================
	% SECTION B.7: EGMONT GROUP AND FIUs
	% ============================================================
	
	\begin{frame}{B.7.1: What Is the Egmont Group and What Does It Do?}
		\begin{block}{Egmont Group - FIU Network}
			\textbf{Answer:} The Egmont Group is a global network of 170+ Financial Intelligence Units.
			\vspace{0.2cm}
			\textbf{Key functions:}
			\begin{itemize}
				\item Secure platform (Egmont Secure Web) for FIU-to-FIU information exchange
				\item Established "Egmont Group Principles for Information Exchange Between FIUs"
				\item Provides training and technical assistance for members
				\item Membership requires meeting criteria (member of FATF or FSRB, operational FIU)
			\end{itemize}
		\end{block}
	\end{frame}
	
	\begin{frame}{B.7.2: What Are the Limitations on Egmont-Exchanged Intelligence?}
		\begin{block}{Egmont Information Exchange Principles}
			\textbf{Answer:} Egmont-exchanged intelligence has strict limitations:
			\vspace{0.2cm}
			\begin{itemize}
				\item Used only for AML/CFT purposes (cannot be used for tax investigations or commercial purposes)
				\item Receiving FIU must maintain confidentiality
				\item Cannot share with third parties without the originating FIU's prior authorization
				\item \textbf{Not automatically admissible in criminal court} (unlike MLAT-obtained evidence)
			\end{itemize}
			\textcolor{red}{Egmont exchanges CANNOT replace MLATs for criminal prosecution.}
		\end{block}
	\end{frame}
	
	% ============================================================
	% SECTION B.8: MUTUAL LEGAL ASSISTANCE TREATIES (MLATs)
	% ============================================================
	
	\begin{frame}{B.8.1: What Is the Purpose of a Mutual Legal Assistance Treaty (MLAT)?}
		\begin{block}{MLAT - Formal Evidence Exchange}
			\textbf{Answer:} MLATs provide formal, legally binding mechanisms for governments to exchange evidence.
			\vspace{0.2cm}
			\textbf{Key features:}
			\begin{itemize}
				\item Exchange of evidence, judicial documents, and witness testimony
				\item Cross-border asset freezing and forfeiture
				\item Records obtained are admissible in criminal proceedings
				\item Each country designates a central authority for processing
			\end{itemize}
			\textcolor{blue}{MLAT vs. Egmont:}
			\begin{itemize}
				\item MLAT: Evidence admissible in court (formal, slow, requires treaty)
				\item Egmont: Intelligence sharing (informal, fast, not court-admissible without conversion)
			\end{itemize}
		\end{block}
	\end{frame}
	
	% ============================================================
	% SECTION B.9: WOLFSBERG GROUP
	% ============================================================
	
	\begin{frame}{B.9.1: What Is the Wolfsberg Group and What Is Its Role?}
		\begin{block}{Wolfsberg Group - Private Banking Association}
			\textbf{Answer:} The Wolfsberg Group is a private association of 13 global banks (Goldman Sachs, Barclays, UBS, Santander, etc.).
			\vspace{0.2cm}
			\textbf{Key functions:}
			\begin{itemize}
				\item Develops best practices and guidance for AML/CFT
				\item Not legally binding, but widely adopted as industry standard
				\item Focus areas: correspondent banking, private banking, AML/CFT, trade finance
				\item Produces the Correspondent Banking Due Diligence Questionnaire (CBDDQ)
			\end{itemize}
			\textbf{Key principle:} Individual risk assessments, not blanket de-risking.
		\end{block}
	\end{frame}
	
	\begin{frame}{B.9.2: What Areas Are Covered in the Wolfsberg CBDDQ?}
		\begin{block}{Wolfsberg CBDDQ - Key Areas}
			\textbf{Answer:} The CBDDQ covers three critical areas:
			\vspace{0.2cm}
			\begin{enumerate}
				\item Respondent bank's AML/CFT policies, procedures, and internal controls
				\item Respondent bank's customer base, including high-risk segments (MSBs, NBFIs, foreign PEPs, shell companies)
				\item Respondent bank's independent testing schedule and results
			\end{enumerate}
			\textbf{Not covered:} CEO's personal investments, office furniture, employee dress code.
			\textcolor{blue}{Most critical risk factor:} Respondent servicing unlicensed MSBs + no recent independent audit.
		\end{block}
	\end{frame}
	
	\begin{frame}{B.9.3: What Is the Wolfsberg Group's Position on De-Risking?}
		\begin{block}{Wolfsberg De-Risking Guidance}
			\textbf{Answer:} The Wolfsberg Group advises against indiscriminate de-risking.
			\vspace{0.2cm}
			\textbf{What is discouraged:}
			\begin{itemize}
				\item Blanket termination of all relationships in a region
				\item Categorical exclusion based solely on geography
				\item No individual risk assessment
			\end{itemize}
			\textbf{What is permissible:}
			\begin{itemize}
				\item Terminating a specific respondent based on documented high-risk factors
				\item Individual, documented risk assessment
			\end{itemize}
			\textcolor{blue}{Justified risk-based exit} (with documentation) is not de-risking; blanket termination is.
		\end{block}
	\end{frame}
	
	% ============================================================
	% SECTION B.10: BASEL COMMITTEE ON BANKING SUPERVISION
	% ============================================================
	
	\begin{frame}{B.10.1: What Does the Basel Committee Expect for AML/CFT Risk Management?}
		\begin{block}{Basel Committee - AML/CFT Integration}
			\textbf{Answer:} The Basel Committee expects AML/CFT risk management to be integrated into the bank's overall operational risk management framework.
			\vspace{0.2cm}
			\textbf{Key expectations:}
			\begin{itemize}
				\item AML/CFT across all three lines of defense (not a siloed compliance function)
				\item Enterprise-Wide Risk Assessments (EWRAs) to identify ML/TF risks
				\item Senior management accountability for AML/CFT
				\item Group-wide AML programs for all branches and subsidiaries
			\end{itemize}
		\end{block}
	\end{frame}
	
	\begin{frame}{B.10.2: What Are the Three Lines of Defense in AML Governance?}
		\begin{block}{Three Lines of Defense Model}
			\textbf{Answer:} The Basel Committee defines three lines of defense:
			\vspace{0.2cm}
			\textbf{First Line - Business Operations:}
			\begin{itemize}
				\item Front-line staff, relationship managers, tellers
				\item Own customer relationships and identify anomalies
			\end{itemize}
			\textbf{Second Line - Compliance:}
			\begin{itemize}
				\item Compliance Department, AML Officer, MLRO
				\item Oversight, guidance, monitoring, reporting
			\end{itemize}
			\textbf{Third Line - Internal Audit:}
			\begin{itemize}
				\item Internal Audit Department
				\item Independent assurance to the Board
			\end{itemize}
			\textcolor{red}{Independence violation:} Audit cannot design controls it later tests.
		\end{block}
	\end{frame}
	
	% ============================================================
	% SECTION B.11: PUBLIC-PRIVATE PARTNERSHIPS (PPPs)
	% ============================================================
	
	\begin{frame}{B.11.1: What Is the Purpose of a Public-Private Partnership (PPP) for AML?}
		\begin{block}{PPP - Intelligence Sharing}
			\textbf{Answer:} PPPs improve intelligence sharing between government authorities and financial institutions.
			\vspace{0.2cm}
			\textbf{Goals:}
			\begin{itemize}
				\item Detect complex financial crime networks (human trafficking, TF cells)
				\item Share emerging typologies and red flags
				\item Help law enforcement prioritize financial intelligence
				\item Improve SAR quality through feedback
			\end{itemize}
			\textbf{Does NOT replace:}
			\begin{itemize}
				\item SAR filing obligations
				\item Independent testing of AML programs
				\item CDD obligations
			\end{itemize}
		\end{block}
	\end{frame}
	
	% ============================================================
	% SECTION B.12: GDPR AND AML CONFLICTS
	% ============================================================
	
	\begin{frame}{B.12.1: How Does GDPR's Right to Erasure Conflict with AML Recordkeeping?}
		\begin{block}{GDPR Article 17(3) - Exception to Right to Erasure}
			\textbf{Answer:} GDPR Article 17(3) explicitly carves out exceptions to the right to erasure where retention is necessary for compliance with a legal obligation.
			\vspace{0.2cm}
			\textbf{Resolution:}
			\begin{itemize}
				\item AML/CFT recordkeeping obligations constitute a legal obligation that overrides the right to erasure
				\item Banks may lawfully refuse to delete records required for AML/CFT (KYC, transaction records, SAR documentation)
				\item Retention is typically 5+ years
				\item Additionally, SAR confidentiality prohibits acknowledging the existence of a SAR to the customer
			\end{itemize}
			\textcolor{red}{AML retention requirements OVERRIDE GDPR erasure requests.}
		\end{block}
	\end{frame}
	
	\begin{frame}{B.12.2: How Should a Bank Handle a U.S. Subpoena for Customer Data from a European Branch?}
		\begin{block}{GDPR vs. U.S. Law Enforcement - Cross-Border Data Transfer}
			\textbf{Answer:} The bank must comply with the U.S. subpoena but can use GDPR compliance mechanisms.
			\vspace{0.2cm}
			\textbf{Resolution:}
			\begin{itemize}
				\item GDPR includes exceptions for cross-border data transfers necessary for important public interest grounds, including law enforcement cooperation
				\item The U.S. branch is subject to U.S. law regardless of GDPR
				\item The bank should seek legal advice on GDPR compliance mechanisms: Standard Contractual Clauses (SCCs), derogations for law enforcement
			\end{itemize}
			\textcolor{blue}{The bank cannot simply ignore the U.S. subpoena because of GDPR.}
		\end{block}
	\end{frame}
	
	% ============================================================
	% SECTION B.13: OFAC EXTRATERRITORIAL REACH
	% ============================================================
	
	\begin{frame}{B.13.1: What Is the Difference Between OFAC Primary and Secondary Sanctions?}
		\begin{block}{Primary vs. Secondary Sanctions}
			\textbf{Answer:} 
			\vspace{0.2cm}
			\textbf{Primary Sanctions:}
			\begin{itemize}
				\item Apply to U.S. persons and U.S.-dollar transactions
				\item Any transaction with a nexus to the U.S. (clearing through a U.S. bank)
				\item Direct jurisdiction - can result in fines and penalties
			\end{itemize}
			\textbf{Secondary Sanctions:}
			\begin{itemize}
				\item Target non-U.S. entities with no U.S. nexus
							\item Conducting "significant" transactions with designated sectors (Iran energy, Russia defense, North Korea)
				\item Consequences: SDN listing, loss of USD correspondent banking access
			\end{itemize}
		\end{block}
	\end{frame}
	
	\begin{frame}{B.13.2: Does the EU Blocking Statute Protect Against U.S. Secondary Sanctions?}
		\begin{block}{EU Blocking Statute vs. U.S. Secondary Sanctions}
			\textbf{Answer:} The EU Blocking Statute provides legal defenses in EU courts but does NOT prevent U.S. secondary sanctions.
			\vspace{0.2cm}
			\textbf{EU Blocking Statute provides:}
			\begin{itemize}
				\item Legal defense in EU courts (prohibits EU companies from complying with U.S. secondary sanctions)
				\item Allows recovery of damages from parties that caused the sanctions
			\end{itemize}
			\textbf{EU Blocking Statute does NOT:}
			\begin{itemize}
				\item Prevent the U.S. from imposing secondary sanctions
				\item Protect USD correspondent banking access
				\item Shield the bank from losing dollar-clearing access
			\end{itemize}
			\textcolor{red}{Practical risk:} U.S. can still cut off the bank's USD correspondent access.
		\end{block}
	\end{frame}
	
	\begin{frame}{B.13.3: How Is the OFAC 50\% Rule Calculated Through Multiple Ownership Tiers?}
		\begin{block}{OFAC 50\% Rule - Multi-Tier Calculation}
			\textbf{Answer:} Ownership is calculated by multiplying through each tier.
			\vspace{0.2cm}
			\textbf{Example:} Individual X (SDN-listed) owns 55\% of Company A. Company A owns 40\% of Company B. Company B owns 60\% of Company C.
			\vspace{0.2cm}
			\textbf{Analysis:}
			\begin{itemize}
				\item Company A is blocked (55\% owned by SDN)
				\item Company B is NOT blocked (40\% ownership from A is below 50\%)
				\item Because Company B is not blocked, its 60\% ownership of Company C does NOT transmit blocked status
				\item Company C is NOT blocked
			\end{itemize}
			\textcolor{red}{If an entity in the chain is not blocked, it does not transmit blocked status downstream.}
		\end{block}
	\end{frame}
	
	% ============================================================
	% SECTION B.14: U.S. BANK SECRECY ACT (BSA)
	% ============================================================
	
	\begin{frame}{B.14.1: What Are the Five Pillars of a BSA/AML Compliance Program?}
		\begin{block}{BSA Five Pillars (31 CFR \S1020.210)}
			\textbf{Answer:} Every financial institution must maintain five minimum pillars:
			\vspace{0.2cm}
			\begin{enumerate}
				\item Internal controls reasonably designed to assure BSA compliance
				\item Designated compliance officer (BSA Officer/MLRO) with authority and resources
				\item Ongoing employee training (initial and periodic, role-based)
				\item Independent testing (audit by independent party, not self)
				\item Customer due diligence (including beneficial ownership identification)
			\end{enumerate}
			\textcolor{red}{All five are required. The CDD Rule (2016) added beneficial ownership as the fifth pillar.}
		\end{block}
	\end{frame}
	
	\begin{frame}{B.14.2: Who Must Be Identified as a Beneficial Owner Under the FinCEN CDD Rule?}
		\begin{block}{FinCEN CDD Rule - Two Prongs of UBO Identification}
			\textbf{Answer:} Two categories of individuals must be identified:
			\vspace{0.2cm}
			\textbf{Ownership Prong:}
			\begin{itemize}
				\item Any individual who directly or indirectly owns 25\% or more of the equity interests of the legal entity
				\item Calculate through tiers (multiply percentages)
			\end{itemize}
			\textbf{Control Prong:}
			\begin{itemize}
				\item At least one individual who exercises significant managerial control
				\item Examples: CEO, CFO, COO, president, general partner, or similar
			\end{itemize}
			\textcolor{red}{The Control Prong is required even if no individual meets the 25\% ownership threshold.}
		\end{block}
	\end{frame}
	
	\begin{frame}{B.14.3: What Protections Does the SAR Safe Harbor Provide?}
		\begin{block}{SAR Safe Harbor (31 U.S.C. \S5318(g))}
			\textbf{Answer:} The Safe Harbor provides specific protections:
			\vspace{0.2cm}
			\textbf{Protected:}
			\begin{itemize}
				\item Immunity from civil liability for good-faith SAR filing
				\item Immunity from liability for failure to disclose the SAR to the customer
			\end{itemize}
			\textbf{Not protected:}
			\begin{itemize}
				\item Using a SAR as evidence in court (except in limited criminal cases for false statements)
				\item Bad-faith filings
			\end{itemize}
			\textbf{Confidentiality:} Unauthorized disclosure of a SAR is a criminal offense (fines + imprisonment up to 5 years). Retention: 5 years from filing date.
		\end{block}
	\end{frame}
	
	\begin{frame}{B.14.4: When Must a Currency Transaction Report (CTR) Be Filed?}
		\begin{block}{CTR Requirements (31 CFR \S1010.311)}
			\textbf{Answer:} CTRs must be filed for cash transactions exceeding \$10,000.
			\vspace{0.2cm}
			\textbf{Key rules:}
			\begin{itemize}
				\item In a single transaction or multiple related transactions
				\item Within a single business day
				\item Filed within 15 days
			\end{itemize}
			\textbf{Exemptions:} Certain customers are exempt (publicly traded companies, government agencies, certain payroll customers).
			\vspace{0.2cm}
			\textcolor{red}{Structuring:} Breaking cash deposits into amounts below \$10,000 to evade CTR is a federal crime (31 U.S.C. \S5324), independent of the underlying source of funds.
		\end{block}
	\end{frame}
	
	% ============================================================
	% SECTION B.15: CORPORATE TRANSPARENCY ACT (CTA)
	% ============================================================
	
	\begin{frame}{B.15.1: What Are the Reporting Requirements Under the Corporate Transparency Act?}
		\begin{block}{CTA - Beneficial Ownership Information (BOI) Reporting}
			\textbf{Answer:} The CTA requires reporting companies to file BOI with FinCEN.
			\vspace{0.2cm}
			\textbf{Key requirements:}
			\begin{itemize}
				\item Reporting companies (LLCs, corporations, other entities) are primarily responsible for filing (not banks)
				\item File initial BOI report within 90 days of formation (30 days for entities created after effective date)
				\item File with FinCEN (not IRS, not state)
				\item Update BOI within 30 days of any change
			\end{itemize}
			\textbf{Exempt entities:} Large operating companies (20+ U.S. employees, \$5M+ revenue, physical U.S. office), publicly traded companies, banks, credit unions.
		\end{block}
	\end{frame}
	
	% ============================================================
	% CONCLUSION SLIDES
	% ============================================================
	
	\begin{frame}{Domain B Summary: Top 20 Most Tested Concepts}
		\begin{enumerate}
			\item \textbf{Mutual Evaluations}: Technical Compliance vs. Effectiveness (11 IOs)
			\item \textbf{Grey List vs. Black List}: Increased Monitoring vs. Call for Action
			\item \textbf{Rec. 10 CDD}: 4 triggers + 4 elements
			\item \textbf{Rec. 16 Travel Rule}: Required originator/beneficiary information
			\item \textbf{Rec. 1 Risk-Based Approach}: Proportionate measures, no de-risking
			\item \textbf{Rec. 3 Predicate Offenses}: All-crimes approach, self-laundering
			\item \textbf{Rec. 5 TF Offense}: No minimum threshold
			\item \textbf{Rec. 19 High-Risk Jurisdictions}: EDD + countermeasures for Black List
			\item \textbf{Rec. 20 STR Filing}: Immediately upon suspicion, no minimum
			\item \textbf{Rec. 21 Tipping-Off}: Prohibited disclosures
			\item \textbf{Rec. 24 Legal Person Transparency}: Bearer shares prohibited or immobilized
			\item \textbf{Rec. 25 Trusts}: Identify settlor, trustee, protector, beneficiaries
			\item \textbf{Rec. 30 Law Enforcement}: Expeditious access to records
			\item \textbf{Rec. 32 Cash Couriers}: Declaration system required
			\item \textbf{USA PATRIOT Sections 311, 313, 314(a), 314(b), 319(a), 319(b)}
			\item \textbf{6AMLD}: 22 predicate offenses, corporate criminal liability
			\item \textbf{5AMLD}: VASPs regulated, lowered prepaid card threshold
			\item \textbf{BSA Five Pillars}: Internal controls, compliance officer, training, testing, CDD
			\item \textbf{Egmont Group}: FIU-to-FIU intelligence sharing (not MLAT replacement)
			\item \textbf{FSRBs}: Autonomous regional bodies, evaluations reported to FATF
		\end{enumerate}
	\end{frame}
	
	\begin{frame}{Critical Red Flags - Domain B Quick Reference}
		\begin{itemize}
			\item \textbf{Low Effectiveness (IO.6)}: FIU receives 60,000 STRs but never disseminates
			\item \textbf{Low Effectiveness (IO.7)}: Zero ML convictions despite high STR volume
			\item \textbf{Low Effectiveness (IO.8)}: Confiscation laws exist but no assets confiscated
			\item \textbf{Rec. 2 Violation}: Supervisor finds AML deficiencies but never tells FIU or LE
			\item \textbf{Rec. 5 Violation}: \$10,000 threshold for TF offense
			\item \textbf{Rec. 24 Violation}: BO registry exists but FIs need court order to access
			\item \textbf{Rec. 30 Violation}: 6-12 month delay to obtain bank records
			\item \textbf{Rec. 32 Violation}: No cash declaration system at borders
			\item \textbf{Section 313 Violation}: U.S. bank maintains correspondent account for shell bank
			\item \textbf{Tipping-Off}: Telling customer "compliance is reviewing your account"
			\item \textbf{De-Risking}: Blanket termination without individual assessment
			\item \textbf{FSRB Evaluation}: Can lead to FATF grey listing
		\end{itemize}
	\end{frame}
	
	\begin{frame}{Exam Day Reminders for Domain B}
		\begin{itemize}
			\item \textbf{Domain B is 20\% of the exam}
			\item FATF Recommendations, Mutual Evaluations, and USA PATRIOT Act are heavily tested
			\item Know the difference between Grey List (EDD required) and Black List (countermeasures required)
			\item SAR filing standard: reasonable suspicion (not conclusive proof)
			\item Distinguish between MLAT (evidence, court-admissible) and Egmont (intelligence, not court-admissible)
			\item FSRBs conduct mutual evaluations using FATF methodology
			\item De-risking (blanket termination) violates the risk-based approach
			\item 6AMLD added corporate criminal liability for legal persons
			\item 5AMLD added VASPs to AML/CFT requirements
			\item CTR threshold in U.S.: \$10,000; SAR: no minimum for ML suspicion
			\item OFAC 50\% Rule: multiply through tiers, only blocked entities transmit blocked status
		\end{itemize}
		\vspace{0.5cm}
		\centering
		\textcolor{blue}{\Large Good luck on your CAMS exam!}
	\end{frame}
	
\end{document}